\documentclass[11pt]{article}
\usepackage{physnotes}

\title{Canonical Quantization of the Scalar Field}
\author{LeeSHEP}
\date{November 2026}
% summary: Real scalar field — Lagrangian, mode expansion, canonical commutation relations.

\begin{document}
\maketitle

\paragraph{Conventions.} Metric signature $(+,-,-,-)$, following Peskin \& Schroeder.

\section{Lagrangian}

The Lagrangian density of a real scalar field is
\begin{equation}
  \Lag = \frac{1}{2}\partial_\mu\phi\,\partial^\mu\phi - \frac{1}{2}m^2\phi^2 .
  \label{eq:lag}
\end{equation}
Applying the Euler--Lagrange equation to \eqref{eq:lag} gives the Klein--Gordon
equation $(\partial^2 + m^2)\phi = 0$.

\section{Mode expansion}

\begin{equation}
  \phi(x) = \int \frac{\dd[3]{p}}{(2\pi)^3}\frac{1}{\sqrt{2E_p}}
  \qty( a_{\mathbf p}\, e^{-\ii p\cdot x} + a^\dagger_{\mathbf p}\, e^{\ii p\cdot x} ),
  \qquad E_p = \sqrt{\mathbf p^2 + m^2}.
\end{equation}
The creation and annihilation operators satisfy
\begin{equation}
  [a_{\mathbf p}, a^\dagger_{\mathbf q}] = (2\pi)^3\,\delta^{(3)}(\mathbf p - \mathbf q).
  \label{eq:ccr}
\end{equation}

\section{Canonical commutator}

Substituting \eqref{eq:ccr} and using $\pi = \dot\phi$,
\begin{equation}
  [\phi(\mathbf x), \pi(\mathbf y)]
  = \int \frac{\dd[3]{p}}{(2\pi)^3}\,\frac{-\ii}{2}
    \qty( -e^{\ii \mathbf p\cdot(\mathbf x-\mathbf y)} - e^{-\ii \mathbf p\cdot(\mathbf x-\mathbf y)} )
  = \ii\,\delta^{(3)}(\mathbf x-\mathbf y).
\end{equation}

\paragraph{Open question.} Beyond normal ordering, how should we understand the
infinite vacuum energy?

\begin{thebibliography}{9}
  \bibitem{ps} M.~E.~Peskin and D.~V.~Schroeder, \emph{An Introduction to Quantum Field Theory}, Ch.~2.
  \bibitem{sred} M.~Srednicki, \emph{Quantum Field Theory}, Ch.~3.
\end{thebibliography}

\end{document}
